\documentclass[a4paper,11pt]{article}
\usepackage[utf8]{inputenc}
\usepackage[T1]{fontenc}
\usepackage[french]{babel}
\usepackage[left=2cm,right=2cm,top=2cm,bottom=2cm]{geometry}
\usepackage{xcolor}
\usepackage{hyperref}
\usepackage{fontawesome5}
\usepackage{parskip}

% --- Couleurs Compagnie des Signaux ---
\definecolor{primary}{RGB}{0, 80, 120}
\hypersetup{colorlinks=true, urlcolor=primary}

\begin{document}

% --- EXPÉDITEUR ---
\noindent
\begin{minipage}[t]{0.5\textwidth}
    \textbf{Loïc Aron MBASSI EWOLO}\\
    Élève-Ingénieur Bac+5 -- Double diplôme ENSEA\\[3pt]
    \faMapMarker\ Cergy, France\\
    \faPhone\ (+33) 7 59 52 33 98\\
    \faEnvelope\ \href{mailto:loicaron.mbassiewolo@ensea.fr}{loicaron.mbassiewolo@ensea.fr}
\end{minipage}
\hfill
% --- DATE ---
\begin{minipage}[t]{0.4\textwidth}
    \raggedleft
    Cergy, le 10 mars 2026
\end{minipage}

\vspace{1.2cm}

% --- DESTINATAIRE ---
\hfill
\begin{minipage}[t]{0.5\textwidth}
    \raggedleft
    \textbf{Compagnie des Signaux}\\
    M. Jean-Pierre Coron\\
    M. Jean-Luc Vendroux\\
    Les Ulis (91), France
\end{minipage}

\vspace{1cm}

\noindent\textbf{Objet :} Candidature en alternance -- Tests composants logiciel embarqué

\vspace{0.8cm}

Madame, Monsieur,

Détecter et isoler des défauts sur un drone en simulant des pannes d'actionneurs et des biais capteurs, puis concevoir un contrôle tolérant aux pannes avec reconfiguration automatique : ce projet académique à l'ENSEA m'a ancré dans la logique des systèmes safety-critical, où tester intégralement le logiciel embarqué n'est pas une option mais une nécessité. C'est cette même exigence que je retrouve dans les calculateurs de sécurité ferroviaire de la Compagnie des Signaux, et qui motive ma candidature.

Mon outil open source Laravel API Generator (+30 étoiles GitHub) illustre concrètement ma capacité à développer, corriger et améliorer un outil logiciel en continu. Ce projet m'a confronté au cycle complet : conception (parsing avancé Python/XML), correction de bugs remontés par les utilisateurs, amélioration itérative des fonctionnalités et rédaction de la documentation technique. Ce workflow correspond aux missions décrites dans votre offre : comprendre un outil de gestion des activités de test, corriger ses bugs, l'améliorer et documenter le tout.

En stage à INRIA Saclay (équipe TRIBE, avril à août 2026), je développe des pipelines Python structurés avec une rigueur de documentation scientifique en anglais. Mon expérience en programmation système C sous Linux (mémoire partagée, signaux POSIX, synchronisation multiprocessus) et en développement embarqué C/C++ sur Nvidia Jetson (challenge CoHoMa, robotique autonome) complète ce profil par une compréhension des contraintes bas niveau et temps réel propres au logiciel embarqué.

Je suis disponible dès septembre 2026 pour une alternance de 12 mois. Intégrer l'équipe OnBoard pour contribuer aux tests des logiciels qui équipent les TGV et les Eurostars représente une opportunité concrète de mettre ma rigueur de programmation et mes compétences Python au service de la sécurité ferroviaire.

Je reste à votre disposition pour un entretien afin de discuter de ma candidature.

Je vous prie d'agréer, Madame, Monsieur, l'expression de mes salutations distinguées.

\vspace{0.8cm}

\textbf{Loïc Aron MBASSI EWOLO}\\
\textit{Élève-Ingénieur ENSEA / Polytechnique Yaoundé}

\end{document}
